% ---- 한글 조판: kotex(xetexko) ----
\usepackage{kotex}
% 이야기 본문은 명조 계열로 통일한다. 메시지·앱 알림·메모는 아래 전용 매크로 글꼴을 유지한다.
% Bundle the OFL fonts used by the book so sandboxed XeTeX renders are stable.
\setmainfont[Path=tex/fonts/]{NanumMyeongjo.ttc}
\setsansfont[Path=tex/fonts/]{NanumGothic.ttc}
\setmonofont[Path=tex/fonts/]{D2Coding.ttc}
\setmainhangulfont[Path=tex/fonts/]{NanumMyeongjo.ttc}
\setsanshangulfont[Path=tex/fonts/]{NanumGothic.ttc}
\setmonohangulfont[Path=tex/fonts/]{D2Coding.ttc}
\newfontfamily\latinx{lmroman10-regular.otf}   % 나눔명조에 없는 라틴 확장 문자(é 등)
\newhangulfontfamily\titlehandwriting{BM Yeonsung} % 배달의민족 연성체: 표제지 전용
% 『시간 고양이』 1권 미리보기 실측 기준: 본문 폭 87mm, 쪽당 19줄, 줄 간격 약 8mm(22.7pt), 글자 약 12~12.5pt, 양끝 맞춤
\makeatletter\renewcommand\normalsize{\@setfontsize\normalsize{12.5pt}{22.7pt}}\makeatother\normalsize
\linespread{1}
\setlength{\parindent}{1em}
\setlength{\parskip}{0pt}
\setlength{\parfillskip}{0pt plus 0.8\linewidth}   % 문단 끝 줄이 너무 짧지 않게(최소 약 1/5)
\setlength{\emergencystretch}{1em}
% ---- 외톨이 줄·쪽 끝 제목 막기 ----
\clubpenalty=10000 \widowpenalty=10000 \displaywidowpenalty=10000
\brokenpenalty=10000
\raggedbottom
\usepackage{needspace}
\usepackage{etoolbox}
\pretocmd{\section}{\Needspace{5\baselineskip}}{}{}
\pretocmd{\subsection}{\Needspace{4\baselineskip}}{}{}
% ---- 머리말·쪽번호 ----
\usepackage[automark]{scrlayer-scrpage}
\automark[chapter]{chapter}
\clearpairofpagestyles
\ofoot*{\pagemark}
% 어린이 독서 흐름을 위해 쪽 위 반복 장 제목(머리글)은 표시하지 않는다.
\ohead{}
\renewcommand*{\chaptermarkformat}{}
\setkomafont{pageheadfoot}{\sffamily\small\color{gray}}
\setkomafont{pagenumber}{\sffamily\small\color{gray}}
\setkomafont{chapter}{\sffamily\bfseries}
\setkomafont{section}{\sffamily\bfseries\large}
\setkomafont{subsection}{\sffamily\bfseries}
\setkomafont{subsubsection}{\sffamily\bfseries}
% ---- 차례: 장 번호 고정 폭 ----
\newcommand{\tocnum}[1]{\makebox[3.3em][l]{#1}}
% ---- 로그 블록(ASCII 화면) ----
\usepackage{fancyvrb}
\usepackage{xcolor}
\usepackage{tikz}
\usetikzlibrary{arrows.meta,positioning,shapes.geometric}
\usepackage{eso-pic}
\definecolor{memoBg}{gray}{0.96}
\definecolor{memoLine}{gray}{0.48}
\definecolor{memoText}{gray}{0.16}
\newcommand{\memobox}[2]{%
  \par\medskip\noindent\fcolorbox{memoLine}{memoBg}{%
    \begin{minipage}{\dimexpr\linewidth-2\fboxsep-2\fboxrule}
    {\sffamily\bfseries\fontsize{11pt}{14pt}\selectfont\color{memoText}#1}\par
    \vspace{1.5mm}
    {\normalfont\fontsize{11.5pt}{18pt}\selectfont\color{black}#2}
    \end{minipage}}\par\vspace{4mm}}
\newcommand{\logfont}{\linespread{1}\fontsize{10.8}{15.2}\selectfont}
% ---- 메신저·앱 알림: 바깥 테두리 없이, 읽기 쉬운 회색 말풍선 ----
\newenvironment{chatthread}{%
  \par\Needspace{5\baselineskip}\medskip\noindent\begin{center}%
  \begin{minipage}{.90\linewidth}\vspace{1.5mm}%
}{%
  \vspace{1.5mm}\end{minipage}\end{center}\par\medskip%
}
\newcommand{\chatroomtitle}[1]{%
  {\noindent\hspace*{\dimexpr.20\linewidth-4.5mm\relax}%
   \sffamily\footnotesize\color{gray}#1\par}\vspace{1mm}%
}
\newcommand{\chatbubble}[2]{%
  \par\noindent\makebox[\linewidth][r]{%
  \begin{tikzpicture}
    \node[anchor=east,rounded corners=3mm,fill=gray!12,inner sep=2.5mm,
      text width=.80\linewidth,align=left] (bubble)
      {{\sffamily\footnotesize\color{gray}#1\par}\vspace{.5mm}%
       \sffamily\fontsize{11.2pt}{15.5pt}\selectfont\color{black}#2};
    \fill[gray!12] ([xshift=-8mm]bubble.south east) --
      ([xshift=-2mm]bubble.south east) --
      ([xshift=-1mm,yshift=-5mm]bubble.south east) -- cycle;
  \end{tikzpicture}}\par\vspace{.5mm}%
}
\newcommand{\appnotice}[2]{%
  \appnoticefrom{#1}{}{#2}%
}
\newcommand{\appnoticefrom}[3]{%
  \par\addvspace{6mm}\Needspace{5\baselineskip}\noindent\makebox[\linewidth][c]{%
  \begin{minipage}{.90\linewidth}
    {\noindent\sffamily\footnotesize\color{gray}#1\par}\vspace{1.5mm}%
    \noindent\colorbox{gray!12}{\parbox{.78\linewidth}{%
      \if\relax\detokenize{#2}\relax\else
        {\sffamily\footnotesize\color{gray}보낸 사람: #2\par}\vspace{1mm}%
      \fi
      \sffamily\fontsize{10.8pt}{15pt}\selectfont\color{black}#3}}%
    \vspace{2mm}%
  \end{minipage}}\par\addvspace{6mm}%
}
\newcommand{\stickynote}[1]{%
  \par\addvspace{7mm}%
  \begingroup\setlength{\fboxsep}{3mm}%
  \noindent\makebox[\linewidth][c]{%
  \fcolorbox{gray!45}{gray!5}{\begin{minipage}{.78\linewidth}
    {\sffamily\fontsize{10.5pt}{16.5pt}\selectfont\color{black}#1\par}
  \end{minipage}}}%
  \endgroup\par\addvspace{8mm}%
}
% ---- 장 첫 쪽 날짜 줄 ----
\newcommand{\dateline}[1]{\par\noindent{\sffamily\small\color{gray}#1}\par\medskip}
% ---- 표제지: 제목·저자·투명 컷아웃 그림을 조판면 중앙에 한 묶음으로 배치 ----
\makeatletter
\renewcommand{\maketitle}{%
  \begin{titlepage}
  \thispagestyle{empty}
  \noindent\vbox to \textheight{%
    \vfil
    \hbox to \textwidth{\hfil
      \begin{minipage}{\textwidth}
      \centering
      {\titlehandwriting\fontsize{19}{26}\selectfont 곰새코지 돌고래 목소리 사건\par}
      \vspace{3mm}
      {\titlehandwriting\fontsize{9.5}{13}\selectfont 이희근\par}
      \vspace{2mm}
      \includegraphics[width=82mm]{images/title-adventure.png}\par
      \end{minipage}%
    \hfil}%
    \vfil
  }%
  \end{titlepage}
}
\makeatother
% ---- 장 도비라 ----
\usepackage{graphicx}
\newcommand{\bodystyle}{\setlength{\parindent}{1em}\setlength{\parskip}{0pt}\normalsize}
\newcommand{\fullpageimage}[1]{%
  \clearpage\thispagestyle{empty}%
  \AddToShipoutPictureBG*{%
    \AtPageLowerLeft{\includegraphics[width=\paperwidth,height=\paperheight]{#1}}%
  }%
  \null\clearpage}
\newcommand{\chapterartpage}[2]{%
  \IfFileExists{#1}{\fullpageimage{#1}}%
    {\IfFileExists{#2}{\fullpageimage{#2}}%
      {\clearpage\thispagestyle{empty}\vspace*{\fill}\begin{center}\fbox{\parbox[c][0.8\textheight][c]{0.8\textwidth}{\centering\color{gray}\small 장 전면 그림이 들어갈 자리}}\end{center}\vspace*{\fill}\clearpage}}}
\newcommand{\dobira}[4]{%
  \clearpage\thispagestyle{plain.scrheadings}\bodystyle%
  \addcontentsline{toc}{chapter}{\tocnum{#1}#2}%
  \markboth{#2}{#2}%
  \vspace*{42mm}%
  {\centering\sffamily{\small #1}\par\vspace{2mm}{\fontsize{19}{27}\selectfont\bfseries #2}\par}%
  \vspace{3mm}%
  \chapterartpage{#4}{#3}}
% ---- 부록(뒷부분): 들여쓰기 없이 문단 사이 띄우기 ----
\newcommand{\backmatterstyle}{\setlength{\parindent}{0pt}\setlength{\parskip}{0.45em plus 0.1em}\backmattersize}
% ---- 판권: 맨 마지막 짝수 쪽 아래 ----
\newcommand{\colophonstart}{\clearpage\thispagestyle{empty}\vspace*{\fill}\begingroup\small\setlength{\parindent}{0pt}\setlength{\parskip}{0.5em}\linespread{1.3}\selectfont}
\newcommand{\colophonend}{\par\endgroup}
% ---- 표: A5 폭 안에 들어오게 ----
\setlength{\tabcolsep}{3pt}
\AtBeginEnvironment{longtable}{\small\linespread{1.2}\selectfont}

% ---- 권말: 본문보다 작은 글자(11pt/18pt) ----
\newcommand{\backmattersize}{\fontsize{10.5pt}{16.5pt}\selectfont}

% ---- 차례 한 쪽에 ----
\RedeclareSectionCommand[tocbeforeskip=1pt plus 1pt]{chapter}
\setkomafont{chapterentry}{\normalfont\sffamily\small}
\setkomafont{chapterentrypagenumber}{\normalfont\sffamily\small}
\AfterTOCHead{\small\linespread{1}\selectfont}
% ---- 앞부분 빈 쪽 줄이기(차례 다음 바로 시작) ----
\AtBeginDocument{\let\cleardoublepage\clearpage}
\RedeclareSectionCommand[beforeskip=6mm,afterskip=5mm]{chapter}
\makeatletter\renewcommand\mainmatter{\clearpage\@mainmattertrue\pagenumbering{arabic}}\makeatother
% ======== 읽기 쉬운 앞부분: 인물 쪽 ========
\definecolor{peopleSea}{HTML}{EAF7F7}
\definecolor{peopleBlue}{HTML}{D9ECFF}
\definecolor{peopleLine}{HTML}{9CB8C9}
\definecolor{peopleText}{HTML}{24404F}
\definecolor{peopleGray}{HTML}{6F7D86}
\newenvironment{peoplepage}{%
  \par\setlength{\parindent}{0pt}\setlength{\parskip}{0pt}%
  \sffamily\fontsize{9.1pt}{12.1pt}\selectfont%
}{\par}
\newcommand{\peopleintro}{%
  \noindent\fcolorbox{peopleLine}{peopleSea}{%
    \begin{minipage}{\dimexpr\linewidth-2\fboxsep-2\fboxrule}
    \sffamily\bfseries\color{peopleText}누가 누구인지 먼저 훑어보기\par
    \vspace{0.2mm}
    {\normalfont\footnotesize\color{peopleGray}아이들 · 가족 · 연구소 · 마을 어른들 · 돌고래가 얽히며 사건이 움직입니다.}
    \end{minipage}}\par\vspace{1.8mm}}
\newcommand{\persongroup}[1]{%
  \par\addvspace{2.1mm}%
  \noindent\colorbox{peopleBlue}{%
    \makebox[\dimexpr\linewidth-2\fboxsep][l]{\sffamily\bfseries\footnotesize\color{peopleText}\quad #1}}%
  \par\vspace{0.8mm}}
\newcommand{\person}[4]{%
  \par\noindent\fcolorbox{peopleLine}{white}{%
    \begin{minipage}{\dimexpr\linewidth-2\fboxsep-2\fboxrule}
    {\sffamily\bfseries\fontsize{10.7pt}{12.5pt}\selectfont\color{peopleText}#1}%
    \if\relax\detokenize{#2}\relax\else\hspace{1.5mm}{\sffamily\scriptsize\color{peopleGray}#2}\fi
    \if\relax\detokenize{#4}\relax\else\hfill{\sffamily\scriptsize\color{peopleGray}소품: #4}\fi\par
    \vspace{0.2mm}
    {\normalfont\fontsize{9.1pt}{11.8pt}\selectfont #3}%
    \end{minipage}}%
  \par\vspace{0.7mm}}
\newenvironment{profilepage}{%
  \par\setlength{\parindent}{0pt}\setlength{\parskip}{0pt}%
  \sffamily\fontsize{8pt}{10pt}\selectfont%
}{\par}
\newcommand{\profileintro}{%
  {\sffamily\footnotesize\color{peopleGray}사진을 먼저 보고, 이름과 역할만 가볍게 훑어보세요.}\par
  \vspace{4mm}}
\newcommand{\profile}[4]{%
  \begin{minipage}[t]{0.305\linewidth}
  \centering
  \fcolorbox{peopleLine}{white}{\includegraphics[width=20mm,height=20mm]{#1}}\par
  \vspace{0.6mm}
  {\sffamily\bfseries\fontsize{10pt}{11.5pt}\selectfont\color{peopleText}#2}\par
  \if\relax\detokenize{#3}\relax\else{\sffamily\scriptsize\color{peopleGray}#3\par}\fi
  \vspace{0.3mm}
  {\normalfont\fontsize{7pt}{8.4pt}\selectfont #4\par}
  \end{minipage}}
\newcommand{\profilecompact}[4]{%
  \begin{minipage}[t]{0.31\linewidth}
  \centering
  \fcolorbox{peopleLine}{white}{\includegraphics[width=25mm,height=25mm]{#1}}\par
  \vspace{0.6mm}
  {\sffamily\bfseries\fontsize{10pt}{11.8pt}\selectfont\color{peopleText}#2}\par
  \if\relax\detokenize{#3}\relax\else{\sffamily\scriptsize\color{peopleGray}#3\par}\fi
  \vspace{0.4mm}
  {\normalfont\fontsize{7.8pt}{9.5pt}\selectfont #4\par}
  \end{minipage}}
\newcommand{\profilehero}[4]{%
  \begin{minipage}[t]{0.31\linewidth}
  \centering
  \fcolorbox{peopleLine}{white}{\includegraphics[width=25mm,height=25mm]{#1}}\par
  \vspace{0.8mm}
  {\sffamily\bfseries\fontsize{10.8pt}{12.5pt}\selectfont\color{peopleText}#2}\par
  \if\relax\detokenize{#3}\relax\else{\sffamily\fontsize{8.8pt}{10.5pt}\selectfont\color{peopleGray}#3\par}\fi
  \vspace{0.7mm}
  {\normalfont\fontsize{7.8pt}{9.8pt}\selectfont #4\par}
  \end{minipage}}
% ======== 이야기 지도형 차례 ========
\usepackage{tocbasic}
\DeclareTOCStyleEntry[
  beforeskip=4pt plus 1.5pt,
  entryformat=\normalfont\sffamily\fontsize{11pt}{19pt}\selectfont,
  pagenumberformat=\normalfont\sffamily\fontsize{10.5pt}{19pt}\selectfont\color{gray},
  linefill=\hfill,
  indent=7mm, numwidth=0pt
]{tocline}{chapter}
\definecolor{tocsea}{HTML}{16879A}
\renewcommand{\tocnum}[1]{\makebox[3.8em][l]{{\color{gray}\footnotesize #1}}}
\newcommand{\tocgroup}[1]{%
  \addvspace{6mm}%
  \noindent\textcolor{tocsea}{\rule{1.2mm}{4.2mm}}\hspace{2.5mm}%
  {\sffamily\bfseries\fontsize{12}{16}\selectfont\color{tocsea}#1}\par
  \addvspace{1.5mm}}
\AfterTOCHead[toc]{\enlargethispage{10mm}}
